% Folder 79, batch 3 (archivists' pages 41-60).
% Pass: Opus 5.5 (claude-opus-5-5), 2026-10-03 — first pass, unchecked against the pages by a human.
%
% Notation fixed for the batch: \Lambda the base ring (Z/nZ), M a free
% Lambda-module of rank 2, P_M the associated "polyhedron" with sommets S_{P_M},
% \pi_0(P_M) its components, P_{M,i} the component i; Gl'_\Lambda(M)^\pm and
% Gl_\Lambda(M)' as he writes them; \mathrm{Autstr} for his "Aut str";
% Z'_\Lambda = \Lambda^*/\{\pm 1_\Lambda\}; \mathbb{S} for his doubled S,
% \mathfrak{P} for his script P (parts), \mathfrak{S}_S for the symmetric group.
\documentclass[11pt,a4paper]{article}
\input{../preamble/grothendieck}

\folder{79}
\batch{3}
\pages{41}{60}
\foldertitle{2-polyèdres réguliers triangulés, et modules libres de rangs 2 sur les anneaux $\Lambda=\mathbb{Z}/n\mathbb{Z}$ et $\Lambda=\mathbb{Z}$ : notes manuscrites (s.d.).}
\dating{[à partir de 1980]}
\watermark{Édition de démonstration}

\begin{document}

\page{41}
\note{feuillet sans numéro de l'auteur. Il continue la démonstration du théorème énoncé au bas de la p.~42 (p.~20 de l'auteur) — « Pour la surjectivité » y fait suite à « L'injectivité … est immédiate » — et sa dernière phrase se poursuit en tête de la p.~43 (p.~21 de l'auteur) : l'ordre de lecture est donc 42, 41, 43.}

Pour la surjectivité, il suffit de prouver celle de
\[
  Gl'_{\Lambda}(M)^{\pm} \longrightarrow \mathrm{Autstr}^{0}(P_M) .
\]
\note{sous « $\mathrm{Autstr}$ », un indice biffé, illisible.}
Soit donc $u\in \mathrm{Autstr}^{0}(P_M)$, choisissons un \struck{\ill{}}
$i \in \pi_0(P_M) \simeq (\det M)^{*}/\pm 1$, $i=\{\pm e_M\}$, et considérons
l'automorphisme de $P_{M,i}$ induit par $u$. Par ce qui \uncertain{précède}, on
\uncertain{a}, $\exists !\, g \in Gl'_{\Lambda}(M)^{\pm}$, tel que
$T_g \mid P_{M,i} = u \mid P_{M,i}$. On est donc ramené au cas où \struck{\ill{}}
$u \mid P_{M,i} = \mathrm{id}_{P_{M,i}}$, et à prouver qu'alors
$u=\mathrm{id}_{P_M}$, ou ce qui revient au même, que $u$ induit l'identité sur
\struck{\ill{}} l'ens des sommets $S_{P_M}$ de $P_M$. Mais cela provient du fait
que $u$ induit l'identité sur \struck{$S_{P_{M,i}}$} $S_{M,i}$, et qu'il respecte
la \struck{\ill{}} structure (101), cqfd.

\emph{Remarque} Dans cette démonstration, on n'a pas utilisé, dans la définition
de $\mathrm{Autstr}(P_M)$, que ces automorphismes respectent la structure \add{(97)}
\uncertain{et} (95), mais seulement (97) et \struck{(1}(101)). \uncertain{Donc} on
voit a posteriori qu'un tel automorphisme
\marginal{en fait, \uncertain{(101)} $\Rightarrow$ (97) et (95), cf. p.~44}
\note{la phrase se poursuit en tête de la p.~43.}

\page{42}
\note{p.~20 de l'auteur. La phrase commence au feuillet précédent de la suite,
hors de ce lot.}

le groupe qu'ils forment, et par $\mathrm{Autstr}^{0}(P_M)$ le sous-groupe formé
des $u$ qui induisent l'identité sur $\pi_0(P_M)$. On trouve donc une suite exacte
canonique
\[
  (102)\qquad 1 \longrightarrow \mathrm{Autstr}^{0}(P_M) \longrightarrow
  \mathrm{Autstr}(P_M) \xrightarrow{\ \delta\ } Z'_{\Lambda} \longrightarrow 1,
\]
où l'hom $\delta$ est défini via l'opération d'un $u\in\mathrm{Autstr}(P_M)$ sur
$\pi_0(P_M)$, qui doit être \struck{\ill{}} de la forme
$i \longmapsto$ \struck{\ill{}} $\delta(u).i$, pour un $\delta(u)\in Z'_{\Lambda}$
bien déterminé (puisque $u$ \struck{conserve} \add{respecte} la structure (97)).
Considérons maintenant l'homom canonique de suites exactes, \struck{déduit}
\struck{de} de la suite exacte
\[
  (103)\qquad 1 \longrightarrow Gl'_{\Lambda}(M)^{\pm} \longrightarrow
  Gl_{\Lambda}(M)' \longrightarrow \Lambda^{*}/\pm 1_{\Lambda} \longrightarrow 1
\]
dans (102), induite par $g \longmapsto T_g$ :

(104)

\begin{tikzcd}[column sep=small]
1 \arrow[r] & Gl'_{\Lambda}(M)^{\pm} \arrow[r] \arrow[d] & Gl_{\Lambda}(M)' \arrow[r] \arrow[d, "g\mapsto T_g"] & \Lambda^{*}/\{\pm 1_{\Lambda}\} \arrow[r] \arrow[d] & 1 \\
1 \arrow[r] & \mathrm{Autstr}^{0}(P_M) \arrow[r] & \mathrm{Autstr}(P_M) \arrow[r] & Z'_{\Lambda} \arrow[r] & 1
\end{tikzcd}

\emph{Théorème} L'hom de suites exactes (104) est un iso.
\note{l'énoncé est marqué d'un trait vertical dans la marge.}

L'injectivité de $g \longmapsto T_g$ est immédiate, elle provient de l'injectivité
\struck{de} de $Gl'_{\Lambda}(M)^{\pm} \longrightarrow \mathrm{Aut}(SP_{(M,e_M)})$,
déjà vue (cette dernière flèche étant en fait bijection).
\note{la démonstration se poursuit p.~41.}

\page{43}
\note{p.~21 de l'auteur ; la première phrase achève celle qui termine la p.~41.}

respecte également la structure (95). Il serait bon d'en comprendre la raison a
priori.

\note{le passage qui suit, jusqu'à « déterminée par la structure (97) », est
encadré et biffé de deux grands traits obliques ; on le transcrit tel qu'il se
lit.}
Pour ceci, on note \add{« (97) (101) »} que l'action d'un automorphisme \struck{de}
$u$ de $P_M$ \add{\uncertain{respectant les structures} \ill{} (101)} est connue,
quand on \struck{la} connaît sur l'ens des sommets $S_{P_M}$, puisque les
\struck{\ill{}} \add{ens des} arêtes et \struck{\ill{}} \add{des} faces
s'identifient à des ens de parties (à deux et à trois éléments respectivement) de
$S_{P_M}$. D'autre part, un automorphisme \add{qui respecte la structure (101)} est
connu quand on connaît son action sur $\pi_0(P_M)$, \emph{et} sur \emph{un}
$S_{P_{M,i}}$. Quand $u$ respecte \add{aussi} la structure (97), l'action de $u$
sur $\pi_0(P_M)$ se décrit via $\delta(u)\in Z'_{\Lambda}$. Lorsque $u$ est de la
forme $T_{\lambda}$, $\lambda\in Z'_{\Lambda}$, on a $\delta(u)=\lambda^{2}$, i.e.
l'action de $T_{\lambda}$ sur $\pi_0(P_M)$ est déterminée par la structure (97).
\marginal{NB On a \ill{} $S_{P_M} \simeq I_M \times \pi_0(P_M)$, \ill{}
$S_{P_M} \supset S_{P_{M,i}}$, idem \ill{} des \ill{}}

Récapitulons le type d'objet structuré qu'on \uncertain{trouve} sur $P_M$.
\struck{Il} \uncertain{On peut le} décrire par les données
\[
  (105)\qquad
  \begin{cases}
    \mathbb{S} \ \text{un ens (de cardinal $\geqslant 3$) (la « valeur commune »}\\
    \qquad \text{des $S_{P_{M,i}}$, pour $i\in\pi_0(P_M)$),}\\
    \pi \ \text{un $Z'_{\Lambda}$-torseur} \quad (Z'_{\Lambda}=\Lambda^{*}/\{\pm 1_{\Lambda}\})\\
    \pi \xrightarrow{\ \Sigma\ } \mathfrak{P}(\mathfrak{P}_3(\mathbb{S})) \ \text{une application}
  \end{cases}
\]
\note{après « un ens », un mot biffé (\uncertain{fini}) ; « (de cardinal
$\geqslant 3$) » est ajouté au-dessus, le 3 corrigé.}
\marginal{en fait card $\mathbb{S}\geqslant 4$}
satisfaisant les conditions nécessaires suivantes :

(106) a) $\forall i\in\pi$, $\Sigma(i)\in\mathfrak{P}(\mathfrak{P}_3(\mathbb{S}))$
est une structure \add{connexe} de 2-polyèdre triangulé sur $\mathbb{S}$,
\struck{\ill{}}

\struck{b) Les $u\in\mathfrak{S}_I$ qui sont des automorphismes}

\page{44}
\struck{(106) b) Soit $G_i = \{u\in\mathfrak{S}_{\mathbb{S}} \mid$ $u$ est un
$\Sigma(i)$-automorphisme de $\mathbb{S}$, $\forall i\in\pi\}$ \ill{}
$\mathrm{Aut}(S,\Sigma(i))$), alors $G_i$ est indépendant de $i$ \ill{}}
\note{deux lignes biffées et barrées de traits obliques.}

(107) b) $\forall i\in\pi$, soit $G_i=\mathrm{Aut}(S,\Sigma(i))$
\add{$\subset\mathfrak{S}_S$}, alors \add{soit $G$,} $G_i$ est indépendant de
$i\in\pi$, (il y a vraiment \ill{} \ill{}), d'autant plus que l'on a c)
ci-dessous)

c) $\pi \xrightarrow{\ \Sigma\ } \mathfrak{P}(\mathfrak{P}_3(S))$ est
\emph{injectif}, i.e. $i\neq j \Longrightarrow \Sigma(i)\neq\Sigma(j)$, plus
précisément,
\[
  i\neq j \Longrightarrow
  \begin{cases}
    A_{\Sigma(i)} \cap A_{\Sigma(j)} = \emptyset\\
    F_{\Sigma(i)} \cap F_{\Sigma(j)} = \emptyset
  \end{cases}
\]
où $A_{\Sigma(i)}$, $F_{\Sigma(i)}$ \add{$=\Sigma(i)$} est l'ens des arêtes resp.
des faces pour \struck{$\Sigma(i)$} la structure $\Sigma(i)$.
\marginal{On a vu \uncertain{aussi} que $G_i^{+}$ est indépendant de $i$, et que
$G$ (resp.\ $G^{+}$) est simplement transitif sur l'ens des régions (resp.\ (108)
des régions directes) des $\Sigma(i)$, i.e.\ que $(S,\Sigma(i))$ est un 2-polyèdre
\emph{régulier} (aux deux sens). \emph{Question} Y a-t-il d'autres structures
triangulées \struck{rég.}\ sur $S$, \uncertain{isomorphes} aux structures \ill{},
et qui ont les mêmes groupes $G$, $G^{+}$ ?}
\note{note écrite en oblique dans la marge gauche, en regard de b) et c).}

Dans les propriétés a) b) c) ci-dessus, la structure de torseur sur $\pi$
n'intervenait pas encore. Introduisant cette structure, on trouve, en désignant
par $A$ le groupe des automorphismes de la structure (105)
\[
  1 \longrightarrow G \longrightarrow A
  \xrightarrow[\text{via action de $A$ sur $\pi$}]{\ \delta\ } Z'_{\Lambda}
\]
et on a

(109) d) L'hom $A \longrightarrow Z'_{\Lambda} = (\mathbb{Z}/n\mathbb{Z})^{*}/\{\pm 1\}$
précédent est surjectif.

On a

(110) e) $\forall \lambda\in Z'_{\Lambda}$, $\exists !\, T_{\lambda}\in\mathfrak{S}_S$,
tel que $\forall s,t\in S$, $i\in\pi$, \struck{\ill{}}
\[
  \bigl[\{s,t\}\in A_{\Sigma(i)} \Longleftrightarrow
  \{T_{\lambda}s, T_{\lambda}t\}\in A_{\Sigma(\lambda i)}\bigr]
\]
L'existence dans les cas qui nous occupe est claire,

\page{45}
\note{p.~22 de l'auteur.}

pour l'unicité, on note que si $u$, $u'$ \struck{\ill{}} $\in S$ \add{et $j\in\pi$}
sont tels que pour tout $s\in S$, on ait
\[
  \bigl(\{s,u\}\in A_{\Sigma(j)} \Longleftrightarrow \{s,u'\}\in A_{\Sigma(j)}\bigr)
\]
\emph{alors} \struck{$z=z'$}. $u=u'$. En d'autres termes,
soit, pour $u\in S$, $i\in\pi$
\[
  (111)\qquad
  \begin{cases}
    A_{u,i} = \text{ens des sommets $\in S$ adjacents à $u$ pour $\Sigma(i)$}
      \quad [\text{card.\ } n]\\
    S_u = \bigcup_{i\in\pi} S_{u,i} \quad
      \text{(réunion disjointe de $\varphi(n)/2$ ens)},
  \end{cases}
\]
\note{le $A$ de $A_{u,i}$ est écrit sur un $S$ ; sous ce symbole, un ajout
interlinéaire, \uncertain{$=\mathrm{Ét}_i(u)$} ; « card.\ $n$ » est corrigé.}
\struck{et soit} des telles façons qu'on a une application de degré $n$
\[
  (112)\qquad S_u \xrightarrow{\ p_u\ } \pi \qquad
  \text{fibre de $i$ est $S_{u,i}$} ;
\]
\note{sous $S_u$, quelques lettres biffées.}
Ceci étant \add{dit,} que si $u,u'\in S$ sont tels que $\forall j\in\pi$, on a
$s\in S_{u,j} \Longleftrightarrow s\in S_{u',j}$ i.e.\ $S_{u,j}=S_{u',j}$,
\struck{donc} alors $u=u'$, on trouve
\[
  (113)\qquad (S_u=S_{u'},\ p_u=p_{u'}) \Longrightarrow u=u',
\]
Mais en fait \struck{$S_u=S_{u'}$, la condition (113) \ill{}}
\struck{\ill{}} signifie aussi que \struck{\ill{}} les formes
$M \longrightarrow \det M$ \struck{\ill{}}
$\varphi_u : z \longmapsto u\wedge z$ et $\varphi_{u'} : z\longmapsto u'\wedge z$
\struck{sont telles}
\struck{\ill{} \ill{} des $z$ tels que $\varphi(z)\in(\det M)^{*}$ \ill{} \ill{}
$\forall x$ tel que $u\wedge x\in(\det M)^{*}$ \ill{} $u\wedge z\equiv 0$ mod $H$
\ill{} $\Longleftrightarrow$ \ill{} $u\wedge(x+z)\in\det M^{*}$}
coïncident \uncertain{dans} l'ens $M^{*}_{u}\wedge M^{*}_{u'}$ des $z\in M$ tels
que $\varphi_u(z)$, $\varphi_{u'}(z)\in(\det M)^{*}$ — d'où on déduit
\add{\ill{}} facilement que \struck{\ill{}} $u'=\pm u$,
\note{les quatre dernières lignes forment un enchevêtrement de biffures, de
cadres et de renvois ; la lecture proposée en suit l'état final et reste
incertaine.}

\page{46}
\note{feuillet sans numéro de l'auteur, entre ses p.~22 et 23.}

(on est ramené à voir que sur un corps \add{\uncertain{$\mathbb{F}_p$}} ayant
\emph{au \struck{moins} trois éléments}, pour un vectoriel $M$ de rg 2 sur ce
corps, le \struck{réunion} complémentaire de la réunion de deux droites engendre
$M$ tout entier …).
\marginal{à expliciter l'argument …}
\note{ce paragraphe est marqué d'un trait vertical dans la marge.}

On a associé au module $M$ libre de rg 2 sur $\Lambda=\Lambda_n=\mathbb{Z}/n\mathbb{Z}$,
le torseur sous $Z'_{\Lambda_n}=Z'_n=(\mathbb{Z}/n\mathbb{Z})^{*}/\pm 1$,
\[
  (114)\qquad \pi = (\det M)^{*}/(\pm 1_{\Lambda}),
\]
à l'aide de celui-ci on reconstitue la « droite mod $\pm 1$ » correspondante
\[
  (115)\qquad L' = (\det M)/\{\pm 1\} \simeq
  (\mathbb{Z}/n\mathbb{Z})/\{\pm 1\} \overset{Z'_n}{\wedge} \pi ;
\]
on a une opération de $Z'_n$ sur $L'$, et $\pi$ en tant que $Z'_n$-ensemble
s'identifie à une orbite de $Z'_n$ dans $L'$ — savoir \add{à} l'unique orbite de
cardinal égal à $\varphi(n)/2=\mathrm{card}\, Z'_n$ :
\[
  (116)\qquad \pi \subset L' .
\]
D'autre part, on associe à $M$ l'ens
\[
  (117)\qquad S = M^{*}/\{\pm \mathrm{id}_M\}
\]
dont on veut saisir les structures essentielles, et les liens entre celles-ci.

\page{47}
\note{p.~23 de l'auteur.}

Un premier élément de structure est une application
\[
  (118)\qquad S\times S \xrightarrow{\ \wedge\ } L' \qquad
  \text{notée $x,y \mapsto x\wedge y$}
\]
\add{(application symétrique) $x\wedge x=0_{L'}$} (obtenue par passage aux
quotients à partir de
\[
  (119)\qquad
  \begin{array}{ccc}
    M^{*}\times M^{*} & \longrightarrow & \det M\\
    (x,y) & \longmapsto & x\wedge y
  \end{array}
  \ \Bigr).
\]
De cette application, on déduit, \struck{\ill{}} pour tout $i\in\pi$
\add{$\subset L'$}, la structure triangulée $\Sigma(i)$ sur $S$, par la formule
\[
  (120)\qquad
  \underbrace{A_{\Sigma(i)}}_{\substack{\text{ens des arêtes pour}\\
  \text{la structure } \Sigma(i)}}
  = \bigl\{\{x,y\}\in\mathfrak{P}_2(S) \mid x\wedge y = i\bigr\}
\]
d'où on déduit l'ens des faces pour la structure $\Sigma(i)$ par
\[
  (121)\qquad F_{\Sigma(i)} = \bigl\{ f\in\mathfrak{P}_3(S) \mid
  a\subset f,\ a\in\mathfrak{P}_2(S) \Longrightarrow a\in A_{\Sigma(i)}\bigr\}
\]
Le fait que cela donne bien une structure triangulée (2-polyèdrale) s'exprime par
les \emph{conditions} \uncertain{propriétés}
\[
  (122)\qquad
  \begin{array}{l}
    \forall x,y\in S \text{ tels que } x\wedge y\in\pi,
      \text{ il existe exactement \emph{deux} éléments $z$ de $S$}\\
    \text{tels que l'on ait } x\wedge y = y\wedge z = z\wedge x .
  \end{array}
\]
\marginal{on voit \uncertain{donc} que dès que card $S\geqslant 2$, \ill{}
card $S\geqslant 4$}
\[
  (123)\qquad \forall x\in S,\ i\in\pi, \text{ soit }
  \begin{cases}
    A_{x,i} = \{y\in S \mid x\wedge y = i\}\\
    F_{x,i} = \bigl\{\{y,z\}\subset S \mid x\wedge y = y\wedge z = z\wedge x = i\bigr\}.
  \end{cases}
\]
Alors $\mathrm{Ét}_{x,i}=(A_{x,i}, F_{x,i}, \text{relation d'incidence})$ est
\add{connexe} (en fait un $n$-polygone ?)
\note{le symbole lu $\mathrm{Ét}_{x,i}$ est surchargé ; lecture incertaine.}

\page{48}
À vrai dire, pour définir les structures triangulées $\Sigma(i)$ ($i\in\pi$), il
suffirait de connaître (118) sur l'ens
\[
  (124)\qquad \underset{\textstyle S\times S}{\overset{\textstyle\cap}{\tilde{A}}}
  = \{(s,t)\in S\times S \mid s\wedge t\in\pi\}
\]
en tant qu'application
\[
  (125)\qquad \tilde{A} \longrightarrow \pi
\]
\struck{\uncertain{appliquée}} \add{$\tilde{A}$,} d'une partie de $S\times S$ dans
$\pi$, \struck{l'image} \add{telle} partie \struck{\ill{}} ne rencontrant pas la
diagonale et est \struck{\ill{}} stable par la symétrie
$(s,t)\mapsto(t,s)$ — donc se donner équivaut à celles d'une partie $A$ de
$\mathfrak{P}_2(S)$, et la donnée de (125), \add{application symétrique,} à
celles d'une application
\[
  (126)\qquad A \xrightarrow{\ t\ } \pi \qquad t=\text{« type »}
\]
dont les fibres sont alors les $A_{\Sigma(i)}$, \add{\ill{}}
\[
  (127)\qquad A = \bigcup_{i\in\pi} \bigl(A_{\Sigma(i)} = t^{-1}(\{i\})\bigr)
\]
On peut dire que la donnée d'une famille, indexée par $\pi$, de structures
triangulées sur $S$, déterminées par les structures de graphes correspondantes
via (121), équivaut à celle d'une partie $A$ de $\mathfrak{P}_2(S)$ (i.e.\ d'une
partie $\tilde{A}$ de $S\times S$ ne rencontrant pas la diag.\ et stable par
symétrie) et d'une application (surjective) (127), de façon que les fibres
$A_i=A_{\Sigma(i)}$

\page{49}
\note{p.~24 de l'auteur ; la p.~48, sans numéro, s'intercale entre ses p.~23 et
24. Ici la numérotation des formules repart à (124), déjà employée p.~48.}

satisfassent (122) et (123). Dans ces données la structure de \add{$Z'_n$-}torseur
sur $\pi$ n'est pas encore intervenue.

\struck{Par contre, quand on fait} \add{Quand on fait} intervenir cette dernière,
$\pi$ se plonge canoniquement dans un $L'$, et la question se pose de décrire,
\emph{en termes des données précédentes}, une application symétrique (118) qui
prolonge (125). La question est liée à celle de décrire une opération de
\struck{$A$} $Z'_n$ sur $S$
\marginal{pas tellement ! \ill{}}
\[
  (124)\qquad \lambda \longmapsto T_{\lambda} : Z'_n \longrightarrow \mathfrak{S}_S
\]
qui satisfasse
\[
  (125)\qquad s\wedge T_{\lambda}t = \lambda(s\wedge t) \qquad \text{pour } s,t\in S,
\]
donc \struck{a fortiori} en particulier
\[
  (125')\qquad
  \begin{array}{ll}
    s\wedge T_{\lambda}t = \lambda(s\wedge t) & \text{pour } (s,t)\in\tilde{A}\\
    \text{i.e.\ } t(\{s,T_{\lambda}t\}) = \lambda . t(\{s,t\}) & \text{pour } \{s,t\}\in A.
  \end{array}
\]
On a vu tantôt qu'une telle \struck{application} \add{opération} \struck{\ill{}}
de $Z'_n$ sur $S$ est uniquement déterminée, \struck{uniquement} du fait que l'on
a
\[
  (126)\qquad
  \begin{cases}
    \text{si $t,t'\in S$ sont tels que } \mathrm{Ét}_{\Sigma(i)}(t) = \mathrm{Ét}_{\Sigma(i)}(t')
      \ \text{pour \struck{tout} } i\in\pi\\
    \qquad \text{(ou seulement pour \emph{un} $i\in\pi$)}\\
    \text{alors } t=t'
  \end{cases}
\]

\page{50}
La relation (125') implique que
\[
  (127)\qquad \{s,t\}\in A \Longrightarrow
  \overbrace{\{T_{\lambda}s, T_{\lambda}t\}}^{=\{s',t'\}}\in A
  \quad\text{et}\quad \tau(\{s',t'\}) = \lambda^{2}\,\tau(\{s,t\}),
\]
donc $T_{\lambda}$ est un automorphisme de la structure
$\{S, \pi \text{ tors.\ sous } Z'_n, A\subset S\times S, \tau\}$,
\struck{\ill{}} \add{ou ce qui revient au même,} de la structure
$\{S, \pi \text{ tors.\ sous } Z'_n, \pi \xrightarrow{\Sigma}$ ens des structures
triangulées sur $S$ comme ens de sommets$\}$, canoniquement associée à ladite
structure.
\note{le type est noté ici $\tau$, là où la p.~48 et la p.~49 écrivent $t$.}
\marginal{NB Il revient au même de se donner l'op.\ de $Z'_n$ sur $\pi$ (en faisant
un $Z'_n$-torseur), ou l'opération $\lambda\mapsto T_{\lambda}$ de $Z'_n$ sur $S$,
reliées par (125') — l'une détermine l'autre.}
\note{note écrite en oblique dans la marge gauche.}

\struck{Remarque : la question de décrire $\wedge$ en termes des autres structures,
on \ill{} pour ceci que l'on a, si}
\[
  \text{\struck{$(128)\quad \{s,t,u\}\in F=\coprod_{i\in\pi} F_{\Sigma(i)},\ v\in S
  \Longrightarrow$}}
\]
\struck{$s\wedge v + t\wedge v + u\wedge v -$ (éventuellement)}
\note{passage barré de grandes croix.}
\marginal{il faut y revenir ultérieurement (118 !)}

J'aimerais comprendre comment \emph{une des structures} $\Sigma(i)$
\emph{détermine toutes les autres}, et l'opération de $Z'_n$ sur l'une de ces
structures.

Soit donc $P_0=(S,A_0,F_0)$ une structure de 2-polyèdre \add{connexe} régulier,
triangulé en l'occurrence, fini par-dessus le marché !
\[
  (128)\qquad \text{Soit}\quad G_0 = \mathrm{Aut}(P_0) \hookrightarrow \mathfrak{S}_S ,
\]

\page{51}
\note{p.~25 de l'auteur ; la p.~50, sans numéro, s'intercale entre ses p.~24 et
25.}

et considérons son \struck{\ill{}} commutant dans $\mathfrak{S}_S$, soit
\[
  (129)\qquad Z = \{\lambda\in\mathfrak{S}_S \mid \lambda u = u\lambda
  \quad \forall u\in G_0\}
\]
Dans le cas où $P_0 = SP_{(M,\{e_M,-e_M\})}$, on a donc
$G_0 \simeq Gl'_{\Lambda}(M)^{\pm}$ opérant sur $S=M^{*}/\{\pm\mathrm{id}_M\}$, son
commutant se réduit à $Z'_n$ (si $\Lambda=\mathbb{Z}/n\mathbb{Z}$).
\marginal{à vérifier}
\struck{Soit $\lambda\in Z$, alors} On a
\[
  (130)\qquad \text{\struck{$\lambda.A_0$}}\ Z\cap G_0 = \mathrm{Cent}(G_0)
  \overset{\mathrm{def}}{=} \mathfrak{z} \subset \mathrm{Cent}(Z)
\]
\note{sous « $\mathrm{Cent}(G_0)$ », un « $=\mathrm{Cent}\,G$ » biffé.}
($\mathfrak{z} = {}_2\struck{Z}Z'$ dans le cas qui nous occupe)
\note{lecture de cette parenthèse incertaine : un symbole biffé, puis
\uncertain{${}_2Z'$}.}
\marginal{donc $\mathfrak{z}=$}
\struck{donc} introduisons aussi
\[
  (131)\qquad Z' = Z/\mathfrak{z}
\]
(on a $Z'=Z$ dans le cas qui nous occupe). Comme la structure \add{$\Sigma_0$}
triangulée sur $S$ est connue quand on connaît $G_0\subset\mathfrak{S}_S$ et
\emph{une} \add{de ses} arêtes, ou \emph{une} de ses faces (puisque $G_0$ est
transitif sur $A_0$ et sur $F_0$), on doit avoir, pour $\lambda\in Z'$
\[
  (132)\qquad A_{\lambda.\Sigma_0}\cap A_{\Sigma_0} = \emptyset, \quad
  F_{\lambda\Sigma_0}\cap F_{\Sigma_0} = \emptyset
\]

\page{52}
plus généralement, \struck{pour $\lambda$, $\lambda'$}
\[
  (132\,\text{bis})\qquad \lambda,\lambda'\in Z',\ \lambda\neq\lambda'
  \Longrightarrow
  \begin{cases}
    A_{\lambda\Sigma_0}\cap A_{\lambda'\Sigma_0} = \emptyset\\
    F_{\lambda\Sigma_0}\cap F_{\lambda'\Sigma_0} = \emptyset .
  \end{cases}
\]
\struck{Donc} D'autre part, \struck{plus généralement}, si
$\lambda\Sigma_0=\Sigma_{\lambda}$, alors on voit que
\[
  (133)\qquad
  \begin{array}{l}
    \lambda \longmapsto \text{\struck{\ill{}}}\ T_{\lambda}(\Sigma_0)
      \overset{\mathrm{def}}{=} \Sigma_{\lambda}\\
    Z' \longrightarrow \text{ens des structures triangulées régulières sur } S
  \end{array}
\]
est une application injective, on a
\marginal{c'est une notation \ill{}, $\Sigma_{\lambda}$ — elle \uncertain{n'a rien
à voir} avec les $\Sigma(i)$ de tout à l'heure, \ill{}
$\Sigma_{\lambda}=\Sigma_{\uncertain{s(\lambda^2)}}$}
\note{note écrite en oblique dans la marge droite, très abrégée.}
\[
  (134)\qquad
  \begin{cases}
    \mathrm{Aut}(S,\Sigma_{\lambda}) = G\\
    \mathrm{Aut}^{+}(S,\Sigma_{\lambda}) = G^{+}
  \end{cases}
  \qquad \text{pour tout } \lambda\in Z' .
\]
Ceci posé, on dira que deux structures 2-polyèdrales triangulées régulières
$\Sigma$, $\Sigma'$ \add{(connexes)} sur un ens.\ de sommets $S$ sont
\emph{équivalentes}, si $\exists\lambda\in\mathfrak{S}_S$, commutant à
$\mathrm{Aut}(S,\Sigma)=G_{\Sigma}$, telle que $\Sigma'=\lambda.(\Sigma)$ — ce qui
implique que
\[
  (135)\qquad
  \begin{cases}
    G_{\Sigma} = G_{\Sigma'}, \quad
      \lambda^{-1}\in\mathrm{Commut}_{\mathfrak{S}_S}(G_{\Sigma'})\\
    \Sigma = \lambda^{-1}.(\Sigma')
  \end{cases}
\]

\page{53}
\note{p.~26 de l'auteur.}

de sorte que la relation envisagée est une relation d'équivalence. Une classe
d'équivalence \add{$\Pi$ pour cette relation} est alors associée à un groupe
d'automorphismes commun $G$ pour les structures envisagées, et se trouve alors
munie de façon naturelle d'une structure de torseur sous le groupe
\[
  (136)\qquad
  \begin{cases}
    Z' = Z/\mathfrak{z}, \quad
      Z \overset{\mathrm{def}}{=} \mathrm{Commut}_{\mathfrak{S}_S}(G)\\
    \mathfrak{z} = Z'\cap G = \mathrm{Cent}(G) .
  \end{cases}
\]
\note{après $Z/\mathfrak{z}$, quelques signes biffés ; « $Z'\cap G$ » est écrit
tel quel, là où (130), p.~51, porte $Z\cap G_0$.}
Le groupe des automorphismes de $(S,\Pi)$ \add{qui induisent l'identité sur $Z'$}
\struck{est alors} s'insère dans une suite exacte

(137)

\begin{tikzcd}[column sep=small]
1 \arrow[r] & G \arrow[r] & \mathrm{Aut}(S,\Pi) \arrow[r] & Z' \arrow[r] & 1 \\
 & \mathfrak{z} \arrow[u, hook] \arrow[r, hook] & Z \arrow[u, hook] \arrow[ur, "\text{épi can.}"'] & &
\end{tikzcd}

Mais cette analyse, dans le cas qui nous occupe, ne nous fournirait pas, à partir
de $SP_{(M,\{\pm e_M\})}=P_{\Sigma_0}$, \struck{toutes} les structures
$SP_{M,i}$, $i\in Z'_n$, mais seulement les structures pour
$i\in(Z'_n)^{2}=\{\lambda^{2} \mid \lambda\in Z'_n\}$, on est encore loin du
compte ! Mais

\page{54}
\struck{notons que pour dans le cas qui nous occupe,}
Mais revenons sur la définition de $Z$ tantôt, en définissant un \add{autre}
groupe, plus petit encore a priori
\[
  (138)\qquad Z_0 = \left\{ \lambda\in\mathfrak{S}_S \;\middle|\;
  \begin{array}{l}
    \lambda \text{ centralise } G_0=\mathrm{Aut}(S,\Sigma_0),\\
    \exists f_{\lambda}\in\mathfrak{S}_S \text{ tel qu'on ait, }
      \forall (s,t)\in S\times S\\
    \{s,t\}\in A_0 \Longleftrightarrow
      \{\underbrace{f_{\lambda}(s)}_{s'}, \underbrace{\lambda^{-1}f_{\lambda}(t)}_{t'}\}\in A_0
  \end{array}
  \right\}
\]
\note{avant « $\forall(s,t)\in S\times S$ », une formule biffée, illisible.}
en d'autres termes, $\lambda\in Z_0$ \struck{ssi pour} $\lambda$ centralise
$G_0=\mathrm{Aut}(S,\Sigma_0)$, et si pour $s',t'\in S$, \struck{la relation} on a
\[
  (139)\qquad \bigl(\{s',\lambda^{-1}t'\}\in A_0 \Longleftrightarrow
  \{\lambda^{-1}s',t'\}\in A_0\bigr) \Longrightarrow (s'\neq t')
\]
de sorte qu'on peut \struck{désigner par $A$} \add{introduire}
\[
  A_{\lambda} = \left\{ \{s',t'\}\in\mathfrak{P}_2(S) \;\middle|\;
  \begin{array}{l}
    \{s',\lambda^{-1}t'\}\in A_0 \text{ i.e.}\\
    \exists (s,t)\in S^{2}, \text{ avec } \{s,t\}\in A_0,\\
    \{s',t'\} = \{s,\lambda.t\}
  \end{array}
  \right\}
\]
et si $A_{\lambda}$ est l'ens des arêtes \struck{pour} d'une structure triangulée sur
$S$ \struck{équivalente} \add{\uncertain{isomorphe}} à la structure $\Sigma_0$
définie par $A_0\subset\mathfrak{P}_2(S)$, i.e.\ s'il existe
$f_{\lambda}\in\mathfrak{S}_S$ telle qu'on ait
\[
  (140)\qquad f_{\lambda}(A_0) = A_{\lambda}
\]
i.e.\ justement tel que
\[
  (141)\qquad \{s,t\}\in A_0 \Longleftrightarrow
  \{f_{\lambda}(s), \lambda^{-1}f_{\lambda}(t)\}\in A_0 .
\]

\page{55}
\note{p.~27 de l'auteur.}

Le fait que $\lambda$ centralise $G_0$ implique que l'on a, pour $u\in G_0$
\[
  (142)\qquad \{s',t'\}\in A_{\lambda} \Longleftrightarrow \{us',ut'\}\in A_{\lambda}
\]
puisque les deux termes de l'équivalence s'écrivent resp.\
\[
  \{s',\lambda^{-1}t'\}\in A_0 \quad\text{et}\quad
  \{us', \underbrace{\lambda^{-1}ut'}_{u\lambda^{-1}t'}\}\in A_0 ,
\]
relations dont l'équivalence signifie justement que $u\in G_0$. \struck{Donc} On
trouve encore
\[
  (143)\qquad \text{\struck{$G_{\lambda}$}}\ \underbrace{\mathrm{Aut}(S,A_{\lambda})}_{G_{\lambda}} = G_0
\]

\emph{Définitions} Soient $S$ un ensemble, $A_0\subset\mathfrak{P}_2(S)$,
$A_1\in\mathfrak{P}_2(S)$ deux structures de graphes \emph{stricts} sur $S$
\add{(i.e.\ une arête comme \uncertain{ensemble} connaît l'ens des \uncertain{ses}
extrémités, et celles-ci \uncertain{sont} distinctes)}, \add{et soient $G_0$ et
$G_1$ les groupes des automorphismes correspondants,} comme ens.\ de ses sommets,
\add{On dit} que ces structures sont \emph{équivalentes}, s'il existe un
$\lambda\in\mathfrak{S}_S$, tel que \struck{l'on} l'on ait
\marginal{On dit alors que $\lambda$ est une transformation de $A_0$ en $A_1$}

a) $\lambda$ centralise $G_0$

b) \struck{\ill{}} \struck{\ill{}} $s,t\in S$, alors \struck{\ill{}}
$\{s,t\}\in A_0 \Longleftrightarrow \{s,\lambda t\}\in A_1$
\note{la définition est marquée d'un trait vertical dans la marge ; $A_1$ est écrit
« $A_1\in\mathfrak{P}_2(S)$ », là où l'on attend $\subset$.}

\emph{Remarque} L'argument précédent montre que l'on a alors
\[
  G_0 = G_1 ,
\]

\page{56}
de sorte que l'on a aussi, pour $s',t'\in S$

a) $\lambda^{-1}\in$ commutant de $G_1$ dans $\mathfrak{S}_S$

b) $\{s',t'\}\in A_1 \Longleftrightarrow \{s',\lambda^{-1}t'\}\in A_0$

— ce qui montre que la relation « d'équivalence » introduite est bien symétrique.
Elle est évidemment réflexive, et la transitivité se voit de même en utilisant que
si $A_0\sim A_1$, $A_1\sim A_2$, \struck{\ill{}} on a
\[
  G_0 = G_1 = G_2
\]
et si
\[
  \begin{array}{l}
    \lambda \text{ transporte } A_0 \text{ en } A_1\\
    \mu \text{ transporte } A_1 \text{ en } A_2
  \end{array}
  \ \Big|\ \lambda\mu \text{ transporte } A_0 \text{ en } A_2
\]
\struck{Notons que} \add{Considérons,} pour $A_0$ fixé, l'\struck{groupe}
\add{ens $\tilde{Z}_0$ des} \struck{les} $\lambda$ tels qu'il \struck{l'on ait}
existe une (évidemment unique) équivalente à $A_0$, $A_1$ et tel que $\lambda$ soit
une transportation de $A_0$ à $A_1$, i.e.\ tels que
\[
  (144)\qquad
  \begin{cases}
    \lambda\in\mathrm{Commut}_{\mathfrak{S}_S}(G_0)\\
    \forall s',t'\in S \quad
      \bigl(\{s',\lambda^{-1}t'\}\in A_0\bigr) \Longleftrightarrow
      \bigl(\{\lambda^{-1}s',t'\}\in A_0\bigr) \Longrightarrow s'\neq t'
  \end{cases}
\]
\marginal{$\forall\{s,t\}\in S$, $\{\lambda s,\lambda^{-1}t\}\in A_0
\Longleftrightarrow \{s,t\}\in A_0$ ; $\{s,t\}\in A_0 \Rightarrow t\neq\lambda s$}
\note{la note marginale, écrite en oblique, est reliée par une flèche à la seconde
ligne de (144) ; sa lecture est incertaine.}
\struck{On aimerait que ce soit un sous-groupe} de $\mathfrak{S}_S$, soit $Z_0$.
\note{le passage « On aimerait que ce soit un sous-groupe » est biffé, les mots
« de $\mathfrak{S}_S$, soit $Z_0$ » ne le sont pas.}

On voit en tous cas que
\[
  (145)\qquad
  \begin{cases}
    \text{\struck{\ill{}}}\quad \mathrm{id}_S\in\tilde{Z}_0\\
    \lambda\in\tilde{Z}_0 \Longrightarrow \lambda^{-1}\in\tilde{Z}_0, \quad
      \lambda,\mu\in\tilde{Z}_0,\ [\lambda,\mu]=1 \Longrightarrow \lambda\mu\in\tilde{Z}_0
  \end{cases}
\]
\note{les tildes sur $Z_0$ sont ajoutés après coup en plusieurs endroits de (145) ;
on les a portés partout où la lecture les suggère.}

\page{57}
\note{p.~28 de l'auteur.}

$\tilde{Z}_0$ est formé d'éléments permutables, p.\ ex.\ si
$\mathrm{Commut}_{\mathfrak{S}_S}(G_0)$ est commutatif, alors $\tilde{Z}_0$ est bien
un sous-groupe de $\mathfrak{S}_S$.
\marginal{NB Si $S\to\mathfrak{P}(S)$, $s\mapsto \mathrm{Ét}(s)$ \ill{} $\Sigma_0$
\uncertain{est injective}, \ill{} $\tilde{Z}_0\ni\lambda,\mu$ \ill{} alors
$\lambda\mu\in\tilde{Z}_0$}
\note{note écrite en oblique dans la marge gauche, en haut de page ; lecture très
incertaine.}

\emph{Définition} On dit que $A_0$ et $A_1$ sont \emph{strictement équivalents},
s'ils sont équivalents et isomorphes, i.e.\ s'il existe \struck{alors}
$\lambda,f\in\mathfrak{S}_S$, tels que

a) $\lambda\in\mathrm{Commut}_{\mathfrak{S}_S}(G_0)$

b) $\forall s,t\in S$, on a
\[
  \{s,t\}\in A_0 \Longleftrightarrow \{f(s),f(t)\}\in A_1
  \Longleftrightarrow \{s,\lambda t\}\in A_1
\]
\note{la première équivalence de b) est surchargée ; lecture incertaine.}
C'est encore une relation d'équivalence.
\struck{Soit} \add{$\lambda\in\tilde{Z}_0$} Soit $Z_0$ \add{l'ens des}
\struck{\ill{}} $\lambda\in\mathfrak{S}_S$ tels que \struck{il existe} le $A_1$
qu'il définit soit \emph{strictement} équivalent à $A_0$, i.e.\ soit isomorphe à
$A_0$. \struck{Alors} sous réserve que $Z_0$ soit formé d'éléments permutables,
$Z_0$ est un sous-groupe de $\mathfrak{S}_S$. On supposera dorénavant qu'il en est
ainsi (i.e.\ que \add{le centralisateur} \add{$\mathrm{Cent}(G_0)$} de $Z_0$ est formé
d'éléments \add{permutables} \ill{} permutables). \struck{Alors l'ens}
\[
  (146)\qquad \mathfrak{z}_0 = Z_0\cap G_0 = \left\{ u\in G_0 \;\middle|\;
  \begin{array}{l}
    u\in\mathrm{Cent}(G_0)\\
    \mathrm{Ét}_{\Sigma_0}(s) = \mathrm{Ét}_{\Sigma_0}(u^{2}s)\\
    \forall s\in S
  \end{array}
  \right\}
\]
\marginal{NB Si $s\mapsto\mathrm{Ét}(s)$, $S\to\mathfrak{P}(S)$ est injective, alors
$\mathfrak{z}=$ \ill{}}
\note{le numéro (146) est mal formé. La fin de la page, entre « d'éléments »
et (146), est surchargée d'ajouts interlinéaires ; la lecture en est incertaine.}

\page{58}
alors l'ens des structures \add{strictement} équivalentes à $A_0$ est un torseur
sous \struck{$Z_0/\mathfrak{z}_0$}
\[
  (147)\qquad Z'_0 \overset{\mathrm{def}}{=} Z_0/\mathfrak{z}_0 .
\]
Donnons-nous, de façon plus intrinsèque, une classe d'équivalence stricte $\Pi$ de
structures de graphes sur $S$, « admissibles » en ce sens que \struck{le commutant
sous-groupe des transformations de $\mathfrak{S}_S$} \add{\ill{}} d'une
$A_0\in\Pi$ vers des autres \struck{\ill{} et \ill{} $A_1\in\Pi$ soit}
$A_1\in\Pi$ soit un groupe, soit $Z_0\subset\mathfrak{S}_S$. Ce groupe ne dépend pas
du choix de $A_0$, pas plus que $G_0=\mathrm{Aut}(S,A_0)$ ne dépend pas du choix de
$A_0\in\Pi$.
\marginal{$\Pi\in\mathfrak{P}(\mathfrak{P}(\mathfrak{P}_2(S)))$, i.e.\
$\Pi\subset\mathfrak{P}(\mathfrak{P}_2(S))$ \ill{}}
\note{la marge gauche porte, en oblique, cette précision sur la nature de $\Pi$, en
partie biffée.}

Soit \struck{Alors}
\[
  (148)\qquad
  \begin{aligned}
    \mathcal{G} = \mathrm{Aut}(S,\Pi) &\overset{\mathrm{def}}{=}
      \bigl\{ f\in\mathfrak{S}_S \mid f(\Pi)=\Pi \bigr\}\\
    &= \bigl\{ f\in\mathfrak{S}_S \mid \forall A_0\in\Pi,\ f(A_0)\in\Pi,\
      f^{-1}(A_0)\in\Pi \bigr\}
  \end{aligned}
\]
\note{la première accolade porte une condition biffée,
\struck{$\forall A_0\in\Pi$, $f(A_0)\in\Pi$ et $f^{-1}(A_0)\in\Pi$}, remplacée
au-dessus par $f(\Pi)=\Pi$ ; la lettre $\mathcal{G}$ est écrite sur un symbole
biffé.}
En termes d'\emph{un} $A_0\in\Pi$ choisi, on a
\[
  (149)\qquad \mathcal{G} = \bigl\{ f\in\mathfrak{S}_S \mid
  f(A_0)\in\Pi,\ f(Z_0)=Z_0 \bigr\} .
\]

\page{59}
\note{p.~29 de l'auteur ; les p.~56 à 58, sans numéro, s'intercalent entre ses
p.~28 et 29.}

On \struck{trouve} s'intéresse au sous-groupe $G$ de $\mathcal{G}$ défini par
\[
  (150)\qquad G = \bigl\{ f\in\mathcal{G} \mid f \text{ centralise } Z_0
  \ \text{\add{mod $\mathfrak{z}_0$,}}\ (\text{i.e.\ } \mathrm{int}(f) \mid Z_0/\mathfrak{z}_0
  = \mathrm{id}_{Z_0/\mathfrak{z}_0}) \bigr\}
\]
Donc on trouve une suite exacte canonique

(151)

\begin{tikzcd}[column sep=small]
1 \arrow[r] & G_0 \arrow[r] & G \arrow[r] & Z'_0 \arrow[r] & 1 \\
1 \arrow[r] & \mathfrak{z}_0 \arrow[r] \arrow[u, hook, "\text{cent}"'] & Z_0 \arrow[r] \arrow[u, hook] & Z_0/\mathfrak{z}_0 \arrow[r] \arrow[u, no head, "\parallel" description] & 1
\end{tikzcd}

\struck{\ill{} que $g$ \ill{} \ill{} si $G$ \ill{} $Z'_0$ \ill{} surjectif}
indépendante du choix d'une « origine » $A_0\in\Pi$.
\note{les flèches verticales de (151) sont dessinées montantes, sans pointe
nette ; on les rend comme des inclusions.}

Supposons maintenant que le groupe $G_0$ soit transitif sur l'ens \add{$A_0$} des
arêtes d'une des structures $A_0\in\Pi$, donc de toutes. On voit alors que
\marginal{$=$ \uncertain{parfaites}}
\[
  (152)\qquad A_0,A_1\in\Pi,\ A_0\neq A_1 \Longrightarrow
  \underbrace{A_0\cap A_1=\emptyset}_{\text{intersection de parties de } \mathfrak{P}_2(S)}
\]
donc si on pose
\[
  (153)\qquad A = A_{\Pi} = \bigcup_{A_i\in\Pi} A_i \subset \mathfrak{P}_2(S)
\]
on a une réunion disjointe donc une application
\[
  (153)\qquad A \xrightarrow{\ \tau\ } \underset{\text{un $Z'_0$-torseur}}{\Pi}
\]
\note{deux formules portent le numéro (153).}

\page{60}
\note{feuillet sans numéro de l'auteur.}

qui généralise (126).

\struck{Quand on part d'une situation \ill{} \ill{}}
Pour appliquer ces réflexions au cas d'un $M$ \struck{libre} de rang 2 sur un
anneau $\Lambda_n=\mathbb{Z}/n\mathbb{Z}$, il faudrait vérifier que pour
$i=\{e,-e\}\in(\det M)^{*}/\{\pm 1_{\Lambda}\}$, d'où un
\[
  A_{M,i} \simeq \mathrm{Rep}(M)/\pm 1 \subset \mathfrak{P}_2(S_M)
\]
\note{entre « $\simeq$ » et « $\subset$ », quelques mots biffés, illisibles.}
où $S_M = M^{*}/\{\pm\mathrm{id}_M\}$, on a \struck{déjà}
\[
  (154)\qquad \underset{\textstyle\mathfrak{S}_S}{\overset{\textstyle\cap}{Z_0}}
  = \text{\struck{\ill{}}}\
  \underbrace{\{T_{\lambda} \mid \lambda\in Z'_n=\Lambda^{*}/\pm 1_{\Lambda}\}}_{\text{noté } (Z'_n)_S}
\]
et on suspecte même que l'on a déjà
\[
  (155)\qquad \mathrm{Commut}_{\mathfrak{S}_S}
  \bigl(\underset{\textstyle Gl'_{\Lambda}(M)^{\pm}}{\overset{\textstyle\Vert}{G_0}}\bigr)
  \overset{?}{=} (Z'_n)_S .
\]
Mais \uncertain{sans s'encombrer} de l'hypothèse $\Lambda\simeq\mathbb{Z}/n\mathbb{Z}$, on peut
regarder sur $S_M = M^{**}/\{\pm\mathrm{id}_M\}$ la structure de graphe
\[
  (i=\{e,-e\}\in(\det M)^{*}/\pm 1_{\Lambda})
\]
définie par $A_{M,i}$ et les groupes $G_0$, $Z_0$ associés.
\note{sous $M^{**}$, une glose : \uncertain{ens des él.\ de $M$ qui font partie
d'une base}.}
\note{l'argument se poursuit au-delà de ce lot.}

\end{document}
